\subsection{Compatibility for connected traffic regions}\label{sec:connectedclassification}
The admissible region family changes the structural answer. Let
\[
\mathcal U^{\mathrm{conn}}_{st}=\{U\subset V:s\in U,\ t\notin U,
\ G[U]\text{ and }G[V\setminus U]\text{ are connected}\},
\]
where connectivity uses the underlying undirected multigraph. These are
separating bonds: both sides of the boundary form connected regions.
\begin{theorem}[two-sided connected-region classification]\label{thm:connectedclassification}
For a connected loopless multigraph, $\mu=e_t-e_s$, and
$D=\operatorname{diag}(w)$, the property
\begin{equation}
\exists\alpha\in\R\quad\forall x\text{ with }Bx=\mu\quad
\forall U\in\mathcal U^{\mathrm{conn}}_{st}:
-\mathbf1_U^\top BDx=\alpha
\label{eq:connectedclassification}
\end{equation}
holds if and only if every edge in the atoms of $\mathcal P(s,t)$ has
one common weight $\alpha$. Weights of edges outside those atoms are
arbitrary, even within an off-path cycle atom.
\end{theorem}
The proof in Appendix~\ref{app:connectedproof} uses connected prefix and
suffix regions within each traversed atom. Each admissible boundary
crosses edges in just one such atom, which explains why off-path
heterogeneity is invisible to this region family. This is a universal
fixed-weight statement; agreement for one realized response need not
imply its conditions.

\begin{corollary}[several queries on connected regions]\label{cor:connectedmulti}
At one shared weight vector, Eq.~\eqref{eq:connectedclassification} holds
for all queried ODs if and only if edges in each group of atoms joined
by the queried block paths have one common weight. An atom on no queried
path has no weight restriction, including no within-atom restriction.
\end{corollary}
Thus restricting attention to connected traffic regions permits additional
outcome heterogeneity. It changes the compatibility classification,
while the unrestricted extrema and error formulas below retain their
own stated boundary family.
